\chapter{Producto Cartesiano}

\section{Definición y ejemplos}

El producto cartesiano entre dos conjuntos es una operación que produce un
nuevo conjunto cuyos elementos son todos los pares ordenados que pueden
formarse, tal que el primer elemento pertenece al primer conjunto y el
segundo elemento pertenece al segundo conjunto.

\begin{definicion}[Par ordenado]
Un par ordenado se compone de dos elementos $x$ e $y$, y se escribe
$(x,y)$, donde $x$ es el primer elemento e $y$ es el segundo elemento. Se
concluye que dos pares ordenados $(x,y)$ y $(z,w)$ son iguales si $x=z$ e
$y=w$.
\end{definicion}

\begin{definicion}[Producto cartesiano]
El \textbf{producto cartesiano} de dos conjuntos $A$ y $B$ es el conjunto,
denotado por $A \times B$, de todos los pares ordenados de la forma
\[
A \times B = \{ (a,b) \mid a \in A \wedge b \in B \}.
\]
\end{definicion}

\begin{ejemplo}
Sean los conjuntos $A = \{1,2,3,4\}$ y $B = \{2,4\}$. El producto cartesiano
$A \times B$ se obtiene colocando al conjunto $A$ como la abscisa y al
conjunto $B$ como la ordenada de un sistema coordenado bidimensional.

\begin{center}
\begin{tikzpicture}[scale=0.8]
  \foreach \x in {1,2,3,4} \draw (\x,0.1) -- (\x,-0.1) node[below] {$\x$};
  \foreach \y in {2,4} \draw (0.1,\y) -- (-0.1,\y) node[left] {$\y$};
  \draw[-Stealth] (0,0) -- (5,0) node[right] {\footnotesize$A$};
  \draw[-Stealth] (0,0) -- (0,5) node[above] {\footnotesize$B$};
  \foreach \x in {1,2,3,4}{
    \foreach \y in {2,4}{
      \fill[colordef] (\x,\y) circle (2pt);
    }
  }
  \node at (4.6,4.6) {\footnotesize $A\times B$};
\end{tikzpicture}
\end{center}

lo que resulta en
\[
A \times B = \{(1,2),(1,4),(2,2),(2,4),(3,2),(3,4),(4,2),(4,4)\}.
\]

Ahora efectuamos el producto cartesiano $B \times A$, colocando esta vez a
$B$ como abscisa y a $A$ como ordenada:
\[
B \times A = \{(2,1),(2,2),(2,3),(2,4),(4,1),(4,2),(4,3),(4,4)\}.
\]
Es claro que $A \times B \neq B \times A$, por lo tanto el producto
cartesiano, en general, \textbf{no conmuta}.
\end{ejemplo}

\section{Cardinalidad del producto cartesiano}

Una pregunta natural es: ¿cuántos elementos posee el conjunto $A \times B$?
En otras palabras, ¿cuál es la cardinalidad de $A \times B$? Sea
$\#(A\times B)$ el número de elementos que contiene $A \times B$.

Sea $A$ un conjunto con $N_A$ elementos, $\#(A) = N_A$, y sea $B$ un
conjunto con $N_B$ elementos, $\#(B) = N_B$. Para encontrar la cardinalidad
de $A \times B$ contamos los elementos que contiene, asumiendo que cada
elemento aparece una sola vez. Tomando el primer elemento de $A$ y formando
todos los pares ordenados posibles con los elementos de $B$, se obtienen
$N_B$ pares ordenados. Tomando el segundo elemento de $A$ se obtienen otros
$N_B$ pares ordenados, y así sucesivamente hasta agotar los $N_A$ elementos
de $A$. Por lo tanto,
\[
\#(A \times B) = \underbrace{N_B + N_B + \dots + N_B}_{N_A} = N_A \cdot N_B .
\]

\begin{proposicionc}
\[
\#(A \times B) = N_A \cdot N_B = \#(A)\cdot\#(B).
\]
\end{proposicionc}

\begin{ejemplo}
Para los conjuntos $A=\{1,2,3,4\}$ y $B=\{2,4\}$ se tiene $\#(A) = 4$ y
$\#(B) = 2$, luego $\#(A\times B) = 4\cdot 2 = 8$.
\end{ejemplo}

\begin{proposicionc}
En general $A \times B \neq B \times A$. Sin embargo, ambos conjuntos
poseen el mismo número de elementos:
\[
\#(A \times B) = N_A \cdot N_B = N_B \cdot N_A = \#(B \times A).
\]
\end{proposicionc}

\subsection{Casos particulares}

Un caso extremo aparece cuando uno de los conjuntos es el conjunto vacío,
$\varnothing$, con $\#(\varnothing) = 0$, por lo que
\[
\#(A \times \varnothing) = N_A \cdot 0 = 0.
\]
El conjunto vacío actúa como el ``cero'' del producto cartesiano. Sin
embargo, si $A$ es un conjunto continuo, entonces $\#(A) = \infty$, y
\[
\#(A \times \varnothing) = \infty \cdot 0 = \; ?
\]
el producto no está determinado por conteo directo. Geométricamente, el
conjunto $A \times \varnothing$ se visualiza como un rectángulo de base $A$
y altura $\varnothing$, es decir, una altura sin elementos; por ello es
razonable que el rectángulo no contenga elementos.

\begin{definicion}
El producto cartesiano de un conjunto $A$ con el conjunto vacío es vacío:
\[
A \times \varnothing = \varnothing .
\]
\end{definicion}

Si uno de los conjuntos posee un solo elemento, por ejemplo $\#(A) = 1$,
entonces
\[
\#(A\times B) = 1\cdot N_B = N_B = \#(B),
\]
por lo tanto, un conjunto con un solo elemento actúa como la \textbf{identidad}
del producto cartesiano.

\begin{axiomac}
Sean $A$ y $B$ dos conjuntos, entonces existe un único conjunto
$A \times B$.
\end{axiomac}

\begin{proposicionc}
Sean $A \neq \varnothing$ y $B \neq \varnothing$, tales que
$A \neq B \Rightarrow A \times B \neq B \times A$.
\end{proposicionc}

\begin{proposicionc}
Sean $A \neq \varnothing$ y $B \neq \varnothing$. Si $A = B \Leftrightarrow
A \times B = B \times A$.
\end{proposicionc}

\begin{proposicionc}
Si $A$, $B$ y $C$ son conjuntos, entonces
\[
(A - B) \times C = (A \times C) - (B \times C).
\]
\end{proposicionc}

\begin{proposicionc}
Si $A \subseteq A'$ y $B \subseteq B'$, entonces $A \times B \subseteq A'
\times B'$.
\end{proposicionc}

\section{Producto cartesiano de conjuntos continuos}

Los casos interesantes aparecen cuando los conjuntos son continuos, de modo
que $\#(A \times B) = \infty$. El producto cartesiano entre conjuntos
continuos genera espacios geométricos interesantes.

\begin{ejemplo}
Sean los intervalos cerrados $A = [1,3]$ y $B = [2,4]$, subconjuntos de
$\mathbb{R}$. El producto cartesiano $A \times B$ es
\[
A \times B = \{ (x,y) \in \mathbb{R}^2 \mid 1 \le x \le 3 \wedge 2 \le y \le 4 \},
\]
que corresponde a un cuadrado compacto en $\mathbb{R}^2$.

\begin{center}
\begin{tikzpicture}[scale=0.9]
  \draw[-Stealth] (0,0) -- (5,0) node[right] {$x$};
  \draw[-Stealth] (0,0) -- (0,5) node[above] {$y$};
  \draw[fill=colordef!25,draw=colordef] (1,2) rectangle (3,4);
  \foreach \v/\l in {1/1,3/3} \draw (\v,0.08) -- (\v,-0.08) node[below] {$\l$};
  \foreach \v/\l in {2/2,4/4} \draw (0.08,\v) -- (-0.08,\v) node[left] {$\l$};
  \node at (2,3) {\footnotesize$A \times B$};
\end{tikzpicture}
\end{center}

Si en cambio se toma $A=[1,3[$ (semiabierto), el conjunto $A$ no contiene a
$x=3$, y esto se refleja en que $A \times B$ no es compacto: el lado
$\{3\} \times [2,4]$ queda excluido de $A \times B$.
\end{ejemplo}

\begin{ejemplo}
Consideremos el producto cartesiano entre un intervalo $A = [a,b] \subset
\mathbb{R}$ y un punto $B = \{c\} \subset \mathbb{R}$:
\[
A \times B = \{(x,y)\in\mathbb{R}^2 \mid a \le x \le b \wedge y = c\}
= \{(x,c)\in \mathbb{R}^2 \mid a\le x\le b\},
\]
que es un segmento de recta horizontal a la altura $y=c$. Análogamente,
$B\times A = \{(c,y)\in\mathbb{R}^2 \mid a\le y \le b\}$ es un segmento
vertical. En general, el producto cartesiano entre un segmento de
$\mathbb{R}$ y un punto genera un segmento de recta en $\mathbb{R}^2$.
\end{ejemplo}

\begin{observacion}
El producto cartesiano entre dos subconjuntos de $\mathbb{R}$ siempre genera
una región rectangular en $\mathbb{R}^2$; por lo tanto, una región de
$\mathbb{R}^2$ que no sea un rectángulo (por ejemplo, el disco unitario
$D = \{(x,y)\in\mathbb{R}^2 \mid x^2+y^2 \le 1\}$) no puede escribirse en la
forma $A \times B$.
\end{observacion}

\section{Producto cartesiano de conjuntos unidimensionales curvos}

El producto cartesiano entre subconjuntos unidimensionales curvos también
da origen a conjuntos en dimensiones mayores. Sea la circunferencia unitaria
\[
S^1 = \{(x,y)\in\mathbb{R}^2 \mid x^2+y^2=1\}
\]
y sea $I = [0,1]$ (usaremos la notación $I$ para un intervalo cerrado de
$\mathbb{R}$ de aquí en adelante).

\begin{ejemplo}
Para encontrar $S^1 \times I$ colocamos $S^1$ como base, horizontalmente en
el plano $Oxy$ de un sistema tridimensional, y en cada punto de $S^1$
colocamos el segmento $I$ en la dirección $z$:
\[
S^1 \times I = \{(x,y,z)\in\mathbb{R}^3 \mid x^2+y^2=1 \wedge 0\le z\le 1\}.
\]
El resultado es el manto de un cilindro, sin las tapas.

\begin{center}
\begin{tikzpicture}[scale=1]
  \draw[colordef, thick] (0,0) ellipse (1.4 and 0.5);
  \draw[colordef, thick] (-1.4,0) -- (-1.4,2.2);
  \draw[colordef, thick] (1.4,0) -- (1.4,2.2);
  \draw[colordef, thick] (0,2.2) ellipse (1.4 and 0.5);
  \node at (2.2,1.1) {\footnotesize$S^1\times I$};
\end{tikzpicture}
\end{center}

También se puede calcular $I \times S^1$, obteniendo
\[
I \times S^1 = \{(x,y,z)\in\mathbb{R}^3 \mid 0\le x\le 1 \wedge y^2+z^2=1\},
\]
nuevamente el manto de un cilindro, pero rotado $90^\circ$ respecto a
$S^1\times I$.
\end{ejemplo}

\begin{ejemplo}[El toro]
El producto cartesiano entre dos circunferencias $S^1$ genera el llamado
\textbf{toro}, denotado
\[
T^2 = S^1 \times S^1 .
\]
Se construye colocando una circunferencia $S^1$ (por ejemplo horizontal) y,
en cada punto de ella, otra circunferencia $S^1$, generando la superficie
de una cámara de bicicleta. Equivalentemente, se puede colocar primero
$S^1$ en forma vertical y en cada punto ubicar otra circunferencia $S^1$;
el resultado es también un toro.

\begin{center}
\begin{tikzpicture}[scale=1.1]
  \draw[colordef, thick] (0,0) ellipse (2.2 and 1.1);
  \draw[colordef, thick] (0,0) ellipse (1.1 and 0.55);
  \draw[colordef, thick] (-2.2,0) .. controls (-2.2,0.55) and (-1.1,0.55) .. (-1.1,0);
  \draw[colordef, thick, dashed] (-2.2,0) .. controls (-2.2,-0.55) and (-1.1,-0.55) .. (-1.1,0);
  \draw[colordef, thick] (2.2,0) .. controls (2.2,0.55) and (1.1,0.55) .. (1.1,0);
  \draw[colordef, thick, dashed] (2.2,0) .. controls (2.2,-0.55) and (1.1,-0.55) .. (1.1,0);
  \node at (0,-1.5) {\footnotesize$T^2 = S^1\times S^1$};
\end{tikzpicture}
\end{center}
\end{ejemplo}

\section{Producto cartesiano en dimensiones mayores}

El producto cartesiano $A\times B$ estudiado hasta ahora se ha realizado con
conjuntos unidimensionales, generando conjuntos bidimensionales. Sin
embargo, también se pueden realizar productos cartesianos entre conjuntos de
mayor dimensión.

\begin{ejemplo}
Sea $A \subset \mathbb{R}^2$ y sea $I=[a,b]\subset\mathbb{R}$. El producto
cartesiano $A \times I$ es
\[
A \times I = \{(x,y,z)\in\mathbb{R}^3 \mid (x,y)\in A \wedge a\le z\le b\},
\]
que genera un cilindro macizo en $\mathbb{R}^3$, es decir, un volumen.
\end{ejemplo}

\begin{observacion}
En el producto $A \times B$, si uno (o ambos) de los conjuntos no es
compacto, entonces $A\times B$ tampoco es compacto.
\end{observacion}

\section{Integración y producto cartesiano}

Dada una función $F(x,y)$ definida en $D \subset \mathbb{R}^2$, el cálculo de
la integral de $F$ sobre $D$ está dado por
\[
I = \int_D F(x,y)\, d^2x, \qquad d^2x = dx\,dy,
\]
donde $D$ se denomina \textbf{región de integración}.

Una región general $D$ se puede escribir en la forma
\[
D := \{(x,y)\in\mathbb{R}^2 \mid x_i \le x \le x_f \wedge f_1(x) \le y \le
f_2(x)\},
\]
de modo que la integral doble se escribe
\[
I = \int_D F(x,y)\, d^2x = \int_{x_i}^{x_f} \int_{f_1(x)}^{f_2(x)}
F(x,y)\, dy\, dx .
\]
La integral se evalúa en dos iteraciones: primero se evalúa la integral
interna (en $y$) y luego la integral externa (en $x$). El proceso no es en
general conmutativo:
\[
\int_{x_i}^{x_f}\int_{f_1(x)}^{f_2(x)} F(x,y)\,dy\,dx \neq
\int_{f_1(x)}^{f_2(x)}\int_{x_i}^{x_f} F(x,y)\,dx\,dy .
\]
Para que el proceso sea conmutativo debemos despejar $y$ en función de $x$ y
reescribir la región $D$ como
\[
D := \{(x,y)\in\mathbb{R}^2 \mid y_i\le y\le y_f \wedge g_1(y)\le x\le
g_2(y)\},
\]
de modo que
\[
I = \int_D F(x,y)\, d^2x = \int_{y_i}^{y_f}\int_{g_1(y)}^{g_2(y)} F(x,y)\,
dx\, dy .
\]
Ambas formas de la integral deben coincidir:
\[
\int_{x_i}^{x_f}\int_{f_1(x)}^{f_2(x)} F(x,y)\,dy\,dx
= \int_{y_i}^{y_f}\int_{g_1(y)}^{g_2(y)} F(x,y)\, dx\, dy .
\]

\begin{ejemplo}[Área del disco compacto]
Sea el disco compacto $D = \{(x,y)\in\mathbb{R}^2 \mid x^2+y^2\le r^2\}$.
Calculemos su área,
\[
A = \int_D dx\,dy .
\]
Reescribiendo $D$ en la forma
\[
D = \{(x,y)\in\mathbb{R}^2 \mid -r\le x\le r \wedge -\sqrt{r^2-x^2}\le y \le
\sqrt{r^2-x^2}\},
\]
la integral queda
\[
A = \int_{-r}^{r} \int_{-\sqrt{r^2-x^2}}^{\sqrt{r^2-x^2}} dy\,dx
= \int_{-r}^{r} y \Big|_{-\sqrt{r^2-x^2}}^{\sqrt{r^2-x^2}} dx
= 2\int_{-r}^{r} \sqrt{r^2-x^2}\, dx .
\]
Con el cambio de variable $x(\alpha) = r\sin\alpha$,
\[
A = 2r^2 \int_{-\pi/2}^{\pi/2} \cos^2\alpha\, d\alpha
= r^2 \left(\alpha + \tfrac12\sin(2\alpha)\right)\Big|_{-\pi/2}^{\pi/2}
\;\Longrightarrow\; A = \pi r^2 .
\]
\end{ejemplo}

\subsection{El caso rectangular \texorpdfstring{$D=I_1\times I_2$}{D=I1xI2}}

El caso de mayor interés es cuando $D = I_1\times I_2$ es un rectángulo en
$\mathbb{R}^2$,
\[
D = I_1\times I_2 = \{(x,y)\in\mathbb{R}^2 \mid x_i\le x\le x_f \wedge
y_i\le y\le y_f\},
\]
de modo que la integral $I$ toma la forma
\[
I = \int_D F(x,y)\, d^2x = \int_{I_1\times I_2} F(x,y)\, d^2x
= \int_{I_1}\int_{I_2} F(x,y)\, dy\, dx
= \int_{x_i}^{x_f}\int_{y_i}^{y_f} F(x,y)\,dy\,dx .
\]
En este caso particular sí se cumple
\[
\int_{I_1\times I_2} F(x,y)\, d^2x = \int_{I_1}\int_{I_2} F(x,y)\,dy\,dx
= \int_{I_2}\int_{I_1} F(x,y)\,dy\,dx = \int_{I_2\times I_1} F(x,y)\, d^2x ,
\]
aun cuando, en general, $I_1\times I_2 \neq I_2\times I_1$ como conjuntos
ordenados.

\begin{observacion}
La integración sobre una región $D=I_1\times I_2$ es distinta de una
integración sobre una región $D = D_1 \cup D_2$, puesto que
\[
I = \int_{I_1\times I_2} F(x,y)\, d^2x = \int_{I_1}\int_{I_2} F(x,y)\,dy\,dx,
\]
\[
I' = \int_{D_1\cup D_2} F(x,y)\, d^2x = \int_{D_1} F(x,y)\,d^2x +
\int_{D_2} F(x,y)\, d^2x,
\]
donde $I_i\subset\mathbb{R}$ y $D_i\subset\mathbb{R}^2$, con $i=1,2$.
\end{observacion}

\subsection{Integrales de una derivada parcial y el borde del rectángulo}

Consideremos el caso $F(x,y) = \dfrac{\partial g(x,y)}{\partial y}$. Entonces
\[
I = \int_D F(x,y)\, d^2x = \int_{I_1\times I_2} F(x,y)\,dy\,dx
= \int_{x_i}^{x_f}\int_{y_i}^{y_f} \frac{\partial g(x,y)}{\partial y}\,
dy\, dx
= \int_{x_i}^{x_f} g(x,y)\Big|_{y=y_i}^{y=y_f} dx,
\]
es decir,
\[
I = \int_{x_i}^{x_f} g(x,y_f)\, dx - \int_{x_i}^{x_f} g(x,y_i)\, dx .
\]

La primera integral está realizada sobre la región $I_1\times\{y_f\}$, que
es uno de los bordes del rectángulo $M = I_1\times I_2$. La segunda integral
está realizada sobre $I_1\times\{y_i\}$, el borde opuesto.

\begin{center}
\begin{tikzpicture}[scale=0.9]
  \draw[-Stealth] (0,0) -- (5,0) node[right] {$x$};
  \draw[-Stealth] (0,0) -- (0,4) node[above] {$y$};
  \draw[fill=colordef!12,draw=colordef] (0.8,0.8) rectangle (3.8,2.8);
  \draw[colorej, ultra thick] (0.8,2.8) -- (3.8,2.8);
  \draw[colorej, ultra thick] (0.8,0.8) -- (3.8,0.8);
  \node at (2.3,3.3) {\footnotesize$I_1\times\{y_f\}$};
  \node at (2.3,0.3) {\footnotesize$I_1\times\{y_i\}$};
  \node at (2.3,1.8) {\footnotesize$M$};
\end{tikzpicture}
\end{center}

Ambas integrales se pueden reescribir en la forma
\[
I = \int_{x_i}^{x_f} g(x,y_f)\,dx - \int_{x_i}^{x_f} g(x,y_i)\,dx
= \int_{I_1} g(x,y_f)\,dx - \int_{I_1} g(x,y_i)\,dx
= \int_{I_1\times\{y_f\}} g(x,y)\,dx - \int_{I_1\times\{y_i\}} g(x,y)\,dx,
\]
donde los puntos $\{y_i\}$ e $\{y_f\}$ actúan aquí como una ``integral
interior''. Esta observación anticipa la relación entre integrales de
derivadas y la frontera de la región de integración, que se desarrolla en el
capítulo siguiente.
